\newcolumntype{Y}{>{\RaggedRight\arraybackslash}X}
\setlist[itemize]{nosep, leftmargin=*, before=\vspace{-0.5\baselineskip}, after=\vspace{-0.5\baselineskip}}

\begin{table}[h!]
    \renewcommand{\arraystretch}{1} % Увеличивает межстрочный интервал
    \setlength{\tabcolsep}{9pt}      % Отступы в ячейках
    \fontsize{10    }{12}\selectfont       % Размер шрифта + интерлиньяж
    \centering
    \caption{Сравнительная таблица методов}
    \label{tab:balancing_schemes}
    \begin{tabularx}{1.1\linewidth}{|X|X|X|}
        \hline
        \fontsize{12}{14}\selectfont\textbf{Схема} & \fontsize{12}{14}\selectfont\textbf{Достоинства} & \fontsize{12}{14}\selectfont\textbf{Недостатки} \\
        \hline

        1. Пасивная балансировка & 
        \begin{itemize}
            \item Простота. Дешевизна. Малый размер.
        \end{itemize} & 
        \begin{itemize}
            \item Не балансируется в простое. 
        \end{itemize} \\
        \hline

        2. Емкостная балансировка & 
        \begin{itemize}
            \item Высокая точность. Простота управления.
        \end{itemize} & 
        \begin{itemize}
            \item Низкая скорость балансировки. Большое количество ключей. Работа только с небольшим напряжением.
        \end{itemize} \\
        \hline
    
        3. Дроссели & 
        \begin{itemize}
            \item Высокая скорость балансировки. Высокий КПД(порядка 95\%)
        \end{itemize} & 
        \begin{itemize}
            \item Сложное управление. 
        \end{itemize} \\
        \hline

        % 4. Трансформаторы & 
        % \begin{itemize}
        %     \item Высокая скорость балансировки.
        % \end{itemize} & 
        % \begin{itemize}
        %     \item Сложное управление.
        % \end{itemize} \\
        % \hline

        5. {\'C}uk конвертеры & 
        \begin{itemize}
            \item Низкие потери. Простота конструкции.
        \end{itemize} & 
        \begin{itemize}
            \item Сложное управление. Необходимость точного измерения напряжений.
        \end{itemize} \\
        \hline

        LAST. ............ & 
        \begin{itemize}
            \item FIXME
        \end{itemize} & 
        \begin{itemize}
            \item FIXME
        \end{itemize} \\
        \hline
    \end{tabularx}
\end{table}

